\documentclass{beamer}
\usepackage[utf8]{inputenc}
\usepackage[french]{babel}
\usepackage{graphics}
\usepackage{tabularx}
\usepackage[french]{babel}


%\usepackage[screen,nopanel]{pdfscreen}
%\usepackage{url}

\usetheme[progressbar=frametitle,numbering=none]{metropolis}
\definecolor{Pourpre}{RGB}{174,37,115} 
\setbeamercolor{alerted text}{fg=Pourpre}
%\setbeamertemplate{itemize items}{\textcolor{Pourpre}{\footnotesize$\blacksquare$}}

\title{Initiation et développement de l'agilité en DUT Informatique}
\author[Thomas Clavier, Yann Secq]
{
  \{ \textbf{Thomas.Clavier \& Yann.Secq} \}\texttt{@univ-lille.fr}
}
\institute{
  \includegraphics[height=13mm]{logo_iut_A_Lille.jpg}
  \includegraphics[height=13mm]{Logo_Univ_de_Lille.jpg}
}
  
\date{}

\logo{
    \includegraphics[width=1cm]{CC_BY_SA} 
}

\begin{document}

\frame{\titlepage}

\begin{frame}{Loi des 2 pieds : bipodocratie}
  \begin{columns}
    \begin{column}{0.5\textwidth}
      \begin{center}
        \includegraphics[width=0.6\textwidth]{foots.png}      
      \end{center}
    \end{column}
    \begin{column}{0.5\textwidth}
      Si vous n’apprenez rien ou que vous ne contribuez pas, passez à autre chose !
    \end{column}
  \end{columns}
\end{frame}

\begin{frame}{2012}
  Journée agile des IUT, Toulouse, Lille, etc. 
\end{frame}

\begin{frame}{2013 : Expérience}
  Premier projet agile, un groupe de S3, 3 jours.

  Objectifs : 
  \begin{itemize}
    \item découvrir l'agilité et l'amélioration continue à travers un vrai projet : radiateur d'information, démo, rétro, PDCA
    \item prototyper rapidement un produit web avec un serveur REST du JS, html et bootstrap.
    \item Vérifier avec un groupe que l'activité est jouable à plus grande échelle. 
  \end{itemize}
\end{frame}

\begin{frame}{2013 : Déroulement}
  \begin{itemize}
    \item 2 encadrants pour un groupe de TP (26 étudiants).
    \item légo4scrum : 2h
    \item les étudiants s'organisent entre eux pour faire des équipes de 6 ou 7.
    \item des sujets apportés par les étudiants
    \item sprint de 2h avec démo et rétro
    \item soutenance en amphie avec démo
    \item rétro de groupe sur l'activité
    \item des contraintes technique forte
  \end{itemize}
\end{frame}

\begin{frame}{2013 : Apprentissages}
  \begin{itemize}
    \item 1 enseignant pour 13 étudiants c'est très bien,
    \item la généralisation à toute la promo est envisageable
    \item avec le budget que nous avons ce sera 5 enseignants pour 4 groupes de TP.
    \item plonger les étudiants dans un contexte technique innovant avec une bonne base théorique (ie. en fin d'IUT) c'est déstabilisant et très formateur. 
    \item activité trop courte pour à la foi faire le backlog, le découper et le coder.
  \end{itemize}
\end{frame}

\begin{frame}{2014 : Éxpérience}
  Fin de S4 décalé, Antoine, 1 porteur de projet externe.

  Toute la promo entre la fin du S3 et le début du S4 : 5 jours du lundi au vendredi, on vole une semaine de congés
  2 équipes sont issue de la journée entrepreneur, 5 enseignants dont 3 non agiliste.
\end{frame}

\begin{frame}{2015 : Éxperiences}
  Début de S3, S4 décallé avec 1 premier porteur de projet, puis toute la promo à la fin du S4, du mercredi au mercredi avec 6 porteurs de projets et les GEA europe, S3 décallé.
\end{frame}

\begin{frame}{Enseignement de l'agilité à l'Université}
Tribulations à l'Université d'un agiliste professionnel et de son compère enseignant-chercheur en quête d'une révolution de la pédagogie académique
   \begin{itemize}
     \item Pourquoi de l'agilité en DUT ?
     \item Comment tout a démarré
     \item Bootstrap itération par itération ...
     \item Et maintenant ?
   \end{itemize}
\end{frame}

\begin{frame}{Pourquoi de l'agilité en DUT}
   \begin{itemize}
    \item DUT $\Rightarrow$ fabrique intensive de développeurs
    \item Contact direct et (assez) étroit avec les réalités professionnelles
    \item Info: algorithmique et programmation, système et réseau, base de données, méthodologie et analyse
    \item Mais enseignement académique $\Rightarrow$ silos disciplinaires :(
    \item MAIS réforme du PPN  $\Rightarrow$ projet de semestre
   \end{itemize}
\end{frame}

\begin{frame}{Comment tout a démarré}
   \begin{itemize}
     \item De l'intérêt d'un réseau national structuré ...
     \item Les "grandes" Journées Agiles(Toulouse'12, Aix'13, Montreuil'14)
     \item Les "petites" Journées Agiles (\textbf{Lille'13}, Lyon'14)
     \item Premier poste de PAST au département pour septembre 2013 !
     \item Début des projets sur la lune pour Laurel et Hardy ;)
   \end{itemize}
\end{frame}

\begin{frame}{Année 1 (septembre 2013 à juillet 2014)}
  \begin{itemize}
    \item de 96h théoriques à XXh pour l'agilité ...
    \item \textbf{Problématique:} \emph{proof of concept}
    \item première expérimentation: à petite échelle
        \begin{itemize}
          \item le groupe "PE" de 26 étudiants
          \item constitution des équipes par les enseignants
          \item 3 jours entre le S3 et le S4 (impact collègues)
          \item sprint de 1h30 + 15mn démo + 15mn rétro
          \item restitution et démo finale \emph{(plus sûr de ça !)}
          \item rétrospective \textbf{très} enthousiaste des étudiants :)
        \end{itemize}
    \item \textbf{Objectif n+1}: généralisation à toute la promotion de S3 !
  \end{itemize}
\end{frame}

\begin{frame}{Année 2 (septembre 2014 à juillet 2015)}
  \begin{itemize}
    \item de 96h théoriques à XXh pour l'agilité ...
    \item \textbf{Problématique:} passage à l'échelle
    \item deuxième expérimentation: large échelle
      \begin{itemize}
        \item ensemble de la promotion (90 étudiants)
        \item besoin d'impliquer des collègues (4 encadrants nécessaires !)
        \item positionnement entre S3 et S4 (pas l'idéal ... mais le plus "pratique")
        \item 5 jours du lundi au vendredi, on "vole" une semaine de congés ;)
        \item 2 équipes issues de la journée entreprenariat sont "protégées"
        \item un \emph{pitch} investisseurs à mi-parcours, un salon à la fin
        \item rétrospective \textbf{très très} enthousiaste des étudiants (mais collègues en difficulté)
      \end{itemize}
      \item \textbf{Objectif n+1}: mini-projet en début de S3 + projet agile final avant départ en stage !
  \end{itemize}
\end{frame}

\begin{frame}{Année 3 (septembre 2015 à juillet 2016)}
  \begin{itemize}
    \item enfin 96h dédiées entièrement à l'agilité :)
    \item \textbf{Problématique:} intégration GEA/INFO + "vrai client" (entrepreneurs) 
    \item troisième expérimentation: l'âge de la maturité ?
      \begin{itemize}
        \item mini-projet agile sur 3j en début de S3
        \begin{itemize}
          \item reboot de POO/COO après les vacances ;)
          \item ensemble de la promotion de S3 (hic: TDD)
          \item collègues de l'année passée poursuivent leur implication :)
          \item un peu trop court, mais première approche
        \end{itemize}
        \item projet agile en S4 décalé
        \begin{itemize}
          \item \emph{comprends pas bien ta phrase ci-dessous: à découper ;)}
          \item Début de S3, S4 décallé avec 1 premier porteur de projet, puis toute la promo à la fin du S4, du mercredi au mercredi avec 6 porteurs de projets et les GEA europe, S3 décallé.
        \end{itemize}
        \item projet agile sur 5j en fin de S4 ...
      \end{itemize}
  \end{itemize}
\end{frame}

\begin{frame}{Année 3 (septembre 2015 à juillet 2016)}
  \begin{itemize}
    \item enfin 96h dédiées entièrement à l'agilité :)
    \item \textbf{Problématique:} intégration GEA/INFO + "vrai client" (entrepreneurs) 
    \item troisième expérimentation: l'âge de la maturité ?
    \begin{itemize}
        \item projet agile sur 5j en fin de S4
        \begin{itemize}
          \item ensemble de la promotion, mais 3 encadrants nécessaire
          \item recrutement d'un agiliste professionnel (plus qu'un naïf dans le coup)
          \item projet arrivant après les cours de REST+JS+Androïd (mieux!)
          \item 5 jours répartis sur 2 semaines (lundi férié, pas de chance)
          \item implication des GEA option "création d'entreprise"
          \item implication de porteur de projets externes: entrepreneurs BGE + Audace
          \item toujours un \emph{pitch} investisseurs à mi-parcours, un salon à la fin
        \end{itemize}
        \item rétrospective \emph{très} enthousiastes des étudiants (effet d'échelle)
      \end{itemize}
      \item \textbf{Objectif n+1}: conserver Laurel (ou Hardy) !
  \end{itemize}
\end{frame}


\begin{frame}{Et maintenant ?}
  \begin{itemize}
    \item Laurel et Hardy, l'aventure continue ?
    \item \'Etape cruciale de la TDD en S3 ...
    \item Renforcement des liens GEA/INFO
      \begin{itemize}
        \item Intégrer les GEA en début de S3 (FI/FC),
        \item équipes communes lors de la journée de l'entreprenariat,
        \item final en beauté avec le projet agile de S4 !
      \end{itemize}
    \item ... et peut-être un jour: DU Lean startup ? :)
  \end{itemize}
\end{frame}

\end{document}

